% Finite forms of Blackwell 1953, Section 5 and Theorems 9--11.
\chapter{Restricted Decisions and Binary Testing}
\label{sec:p17-restricted-decisions-binary-testing}

A universal information comparison concerns every finite decision problem.
Sometimes the number of available decisions is fixed. This chapter develops
that restriction, proves that two decisions suffice when there are two
states, and derives the testing-error comparison and its preservation under
repeated conditionally independent observations.
The source is Blackwell's Section~5 and Theorems~9--11
\cite[pp.~269--272]{Blackwell1953EquivalentExperiments}.
All experiments here have finite signal sets. Their definitions, value
function, garbling relation and full comparison theorem are those of
\cref{sec:p17-value-of-information}.

\section{A fixed number of decisions}
\label{sec:p17-k-decision-comparison}

\begin{hbdefinition}{def:p17-k-output-experiments}{Experiments obtainable with a fixed output alphabet}
Fix \(k\geq1\). For an \(n\times r\) experiment \(P\), define
\[
 \mathcal P_{P,k}
   =\{PM:M\text{ is a row-stochastic }r\times k\text{ matrix}\}.
\]
For \(f:\{1,\ldots,r\}\to\{1,\ldots,k\}\), put \(M^f_{zj}=1\)
if \(f(z)=j\), and zero otherwise, and define
\[
                 \mathcal P^*_{P,k}=\{PM^f:f\text{ is such a map}\}.
\]
The sets \(f^{-1}(j)\) form a labelled partition; empty cells are allowed.
Thus the starred set uses deterministic processing and the unstarred set
permits randomized processing. For experiments on the same states, write
\[
              P>_kQ\quad\Longleftrightarrow\quad
                    \mathcal P_{P,k}\supseteq\mathcal P_{Q,k}.
\]
This is the source's non-strict \(k\)-decision comparison.
An action list of length at most \(k\) can be padded by repeated actions,
so at most \(k\) actions and a \(k\)-entry list give the same comparison.
\index{information comparison!restricted decisions}
\end{hbdefinition}

\begin{hblemma}{thm:p17-randomized-decision-hulls}{proved=annex:p17.randomized-decision-hulls}{Finite randomization and convex hulls}
For every \(P,k\),
\[
           \mathcal P_{P,k}=\operatorname{conv}\mathcal P^*_{P,k}.
\]
If \(C\) is any nonempty set of statewise loss vectors and
\(\mathcal A=\operatorname{conv}C\), then
\[
                    B(P,\mathcal A)=\operatorname{conv}B(P,C),
\]
where \(B(P,C)\) consists of the loss vectors obtained by selecting an
element of \(C\) at each signal. If
\(\mathcal A=\operatorname{conv}\{a_1,\ldots,a_k\}\) and \(D\) has
these vectors as rows, then
\[
           B(P,\mathcal A)
                  =\{\operatorname{diag}(RD):R\in\mathcal P_{P,k}\}.
\]
The proof gives explicit product weights for all deterministic selections.
No limiting purification argument is needed for a finite signal alphabet.
\end{hblemma}

\begin{hbtheorem}{thm:p17-finite-k-comparison}{proved=annex:p17.finite-k-comparison}{Equivalent finite restricted-decision comparisons}
For \(k\geq1\) and experiments \(P,Q\) on the same states, the following
conditions are equivalent.
\begin{enumerate}
\item \(P>_kQ\).
\item \(B(P,\mathcal A)\supseteq B(Q,\mathcal A)\) for every
      \(\mathcal A\) that is the convex hull of at most \(k\) loss vectors.
\item \(V(\pi,P,u)\geq V(\pi,Q,u)\) for every prior and every nonempty
      action set of size at most \(k\), with arbitrary real payoffs.
\item The inequality in item 3 holds for one specified full-support
      prior and every such action/payoff problem.
\item At that same full-support prior \(\pi\),
      \[
       \int\phi\,d\mu_{P,\pi}\geq\int\phi\,d\mu_{Q,\pi}
      \]
      for every function \(\phi\) on \(\Delta(S)\) that is a maximum of
      at most \(k\) linear functions.
\end{enumerate}
The integrals denote the finite posterior sums already defined.
The source uses the uniform prior in its standard-measure formulation;
the fixed-full-support-prior formulation above is proved explicitly.
\end{hbtheorem}

\begin{hbcorollary}{thm:p17-k-comparison-hierarchy}{proved=annex:p17.k-comparison-hierarchy}{Hierarchy and the finite output bound}
The relation \(>_k\) is reflexive and transitive. Every pair of experiments
is mutually comparable under \(>_1\). For every \(k\geq1\),
\[
              P>Q\ \Longrightarrow\ P>_{k+1}Q\
                         \Longrightarrow\ P>_kQ.
\]
If \(Q\) has \(t\) signals and \(k\geq t\), then
\[
                         P>_kQ\quad\Longleftrightarrow\quad P>Q.
\]
Dominance for all finite decision bounds is therefore equivalent to full
garbling dominance here. The bound depends on the compared experiment's
number of signals; no equivalence of every smaller decision bound in
every state dimension is asserted.
\end{hbcorollary}

\section{Why two decisions suffice for two states}
\label{sec:p17-two-state-comparison}

\begin{hblemma}{thm:p17-convex-hinge-decomposition}{proved=annex:p17.convex-hinge-decomposition}{Convex polygonal functions on the binary simplex}
Identify the two-state simplex with \(x\in[0,1]\) by
\((p_1,p_2)=(x,1-x)\). Any maximum of finitely many linear functions
on this simplex has a representation
\[
       \phi(x)=a+bx+\sum_{j=1}^m d_j(x-t_j)_+,
       \qquad
       0<t_1<\cdots<t_m<1,\quad d_j\geq0,
\]
where \(v_+=\max(v,0)\). The empty sum is allowed.
Each hinge is the maximum of two linear functions on the simplex:
\[
        (x-t)_+=\max\{0,(1-t)p_1-tp_2\}.
\]
The affine part also restricts a linear function of \((p_1,p_2)\).
The annex constructs the breakpoints and coefficients and proves the
nonnegative slope increments.
\end{hblemma}

\begin{hbtheorem}{thm:p17-binary-two-decisions}{proved=annex:p17.binary-two-decisions}{Two-state completeness of two-decision comparison}
For two-state finite experiments,
\[
                              P>_2Q\quad\Longleftrightarrow\quad P>Q.
\]
Equivalently, it suffices to compare all two-action payoff problems at
one fixed full-support prior. This is the finite form of source
Theorem~10. The proof supplies its polygonal convex-function decomposition
\cite[pp.~271--272]{Blackwell1953EquivalentExperiments}.
\end{hbtheorem}

\section{The complete attainable testing-error region}
\label{sec:p17-binary-testing-region}

\begin{hbdefinition}{def:p17-binary-test}{Randomized tests and the two errors}
State 1 is the null hypothesis and state 2 the alternative. For \(P\),
a test is a vector \(\delta=(\delta_z)\in[0,1]^r\).
After signal \(z\), it rejects the null with probability \(\delta_z\).
Its error probabilities are
\[
    a_P(\delta)=\sum_zp_{1z}\delta_z,\qquad
    b_P(\delta)=\sum_zp_{2z}(1-\delta_z).
\]
These are the error of the first kind and of the second kind, respectively.
Power means \(1-b_P(\delta)\). Define
\[
 \mathcal R(P)=\{(a_P(\delta),b_P(\delta)):\delta\in[0,1]^r\},
 \qquad
 f_P(t)=\min\{b:(t,b)\in\mathcal R(P)\}\quad(0\leq t\leq1).
\]
This constraint fixes the first error exactly. The constant test
\(\delta_z=t\) makes every section nonempty; the next theorem proves its
minimum exists and relates it to an upper-bound constraint.
\index{hypothesis testing!attainable error region}
\index{receiver operating characteristic}
\end{hbdefinition}

\begin{hbtheorem}{thm:p17-testing-error-region}{proved=annex:p17.testing-error-region}{Geometry and ordering of the testing-error curves}
The region \(\mathcal R(P)\) is a compact convex polygon, possibly a
line segment. It contains \((0,1),(1,0)\) and is invariant under
\((a,b)\mapsto(1-a,1-b)\). Its full description is
\[
 \mathcal R(P)
   =\{(t,b):0\leq t\leq1,\quad f_P(t)\leq b\leq1-f_P(1-t)\}.
\]
The function \(f_P\) is convex and nonincreasing, with \(f_P(1)=0\), and
\[
 f_P(t)=\min_{\delta:a_P(\delta)\leq t}b_P(\delta).
\]
For two-state finite experiments,
\[
 \begin{aligned}
 P>Q
 &\quad\Longleftrightarrow\quad \mathcal R(P)\supseteq\mathcal R(Q)\\
 &\quad\Longleftrightarrow\quad
                   f_P(t)\leq f_Q(t)\quad(0\leq t\leq1).
 \end{aligned}
\]
The entire optimal testing-error curve characterizes the information
comparison. Equality of exact-size and size-at-most optima uses randomized
tests; it need not hold when tests must be deterministic.
\end{hbtheorem}

\begin{hbmethod}{def:p17-finite-likelihood-test}{Constructing the finite optimal curve}
Write \(p_z=p_{1z}\), \(q_z=p_{2z}\). First reject at every signal
with \(p_z=0<q_z\); these have no first-kind cost. Signals with both
masses zero have no effect. For the remaining signals \(p_z>0\), order
\(\ell_z=q_z/p_z\) from largest to smallest, breaking ties arbitrarily.
Put
\[
 A_j=\sum_{i=1}^j p_{z_i},\qquad
 B_j=\sum_{i=1}^j q_{z_i},\qquad
 q_*=\sum_{z:p_z=0}q_z,\qquad A_0=B_0=0.
\]
For \(0\leq t<1\), take the first \(j\) with \(A_j>t\).
Reject at preceding positive-\(p\) signals, randomize at \(z_j\) with
probability \((t-A_{j-1})/p_{z_j}\), and accept at later signals.
At \(t=1\), reject everywhere. The resulting optimal curve is
\[
      f_P(t)=1-q_*-B_{j-1}
                  -\ell_{z_j}(t-A_{j-1})\quad(0\leq t<1).
\]
The construction includes ties, zeros and boundary sizes.
It is a finite likelihood-ratio test. Plotting power \(1-f_P(t)\)
against allowed first-kind error \(t\) gives its piecewise-linear
receiver operating characteristic.
\end{hbmethod}

\begin{hbtheorem}{thm:p17-finite-likelihood-optimality}{proved=annex:p17.finite-likelihood-optimality}{Optimality of the ratio construction}
The procedure in \cref{def:p17-finite-likelihood-test} attains \(f_P(t)\)
at every \(t\). Any ordering within a tied ratio block gives the same
optimal curve. In particular,
\[
             f_P(0)=1-\sum_{z:p_{1z}=0}p_{2z}.
\]
The lower boundary need not pass through \((0,1)\): positive alternative
mass on null-impossible signals gives power at size zero.
The entire error region still contains the always-accept endpoint.
\end{hbtheorem}

\section{Independent observations and worked comparisons}
\label{sec:p17-repeated-information}

\begin{hbdefinition}{def:p17-repeated-experiment}{A conditionally independent repeated experiment}
For an integer \(m\geq1\), define \(P^{[m]}\) on signal tuples
\(z=(z_1,\ldots,z_m)\) by
\[
                         P^{[m]}_{s,z}=\prod_{i=1}^m p_{s,z_i}.
\]
The state \(s\) is the same throughout the sample. Observations are
independent conditional on that state, with the same conditional law.
Under a mixture over states they need not be independent without conditioning.
\end{hbdefinition}

\begin{hbtheorem}{thm:p17-repeated-experiment-comparison}{proved=annex:p17.repeated-experiment-comparison}{Product garbling and preservation for every sample size}
If \(PM=Q\), the row-stochastic matrix
\[
                M^{[m]}_{z,y}=\prod_{i=1}^m M_{z_i,y_i}
\]
satisfies \(P^{[m]}M^{[m]}=Q^{[m]}\).
All finite-state decision comparisons are therefore preserved by the
conditionally independent repetition defined above.
For two states, if \(f_P(t)\leq f_Q(t)\) at every size with one observation,
then
\[
                     f_{P^{[m]}}(t)\leq f_{Q^{[m]}}(t)
                    \quad(0\leq t\leq1,\ m\geq1).
\]
This supplies the finite product proof for source Theorem~11
\cite[p.~272]{Blackwell1953EquivalentExperiments}.
\end{hbtheorem}

\begin{hbexample}{def:p17-binary-testing-calculation}{The full curve for a binary symmetric signal}
For
\[
                  P=\begin{pmatrix}4/5&1/5\\1/5&4/5\end{pmatrix},
\]
the second signal has likelihood ratio \(4\), and the first has ratio
\(1/4\). Therefore
\[
 f_P(t)=
 \begin{cases}
  1-4t,&0\leq t\leq1/5,\\
  (1-t)/4,&1/5\leq t\leq1.
 \end{cases}
\]
For the earlier garbling
\(Q=\bigl(\begin{smallmatrix}13/20&7/20\\7/20&13/20\end{smallmatrix}\bigr)\),
\[
 f_Q(t)=
 \begin{cases}
  1-(13/7)t,&0\leq t\leq7/20,\\
  (7/13)(1-t),&7/20\leq t\leq1.
 \end{cases}
\]
At size \(1/5\), their errors are \(1/5\) and \(22/35\).
For two independent \(P\)-observations, the tuple ratios are
\(16,1,1,1/16\), and
\[
 f_{P^{[2]}}(t)=
 \begin{cases}
  1-16t,&0\leq t\leq1/25,\\
  2/5-t,&1/25\leq t\leq9/25,\\
  (1-t)/16,&9/25\leq t\leq1.
 \end{cases}
\]
At size \(1/10\), one observation has error \(3/5\) and two have error
\(3/10\). The tuple product assumption is part of this calculation.
\end{hbexample}

\begin{hbtheorem}{thm:p17-testing-calculations-boundaries}{proved=annex:p17.testing-calculations-boundaries}{Testing calculations, deterministic sizes and dependent repetitions}
The preceding curves and values hold. In addition:
\begin{enumerate}
\item A single-signal experiment permits deterministic first-kind errors
      only 0 and 1; randomized tests permit every size.
\item For mutually disjoint point-mass null and alternative laws,
      \(f_P(t)=0\) for every \(t\), including zero.
\item The product theorem needs its independence premise. There are two
      joint experiments whose one-observation marginals coincide and are
      uninformative, but whose joint information values differ, with no
      garbling from the uninformative joint experiment to the informative one.
\end{enumerate}
The annex supplies all outcome tables and an explicit distinguishing decision.
\end{hbtheorem}

\begin{proofboundary}{def:p17-restricted-testing-scope}{Finite testing comparison and further regimes}
The constructions cover all finite action bounds and sample sizes for
finite signal alphabets. The general observation-space source theorems
require measure-theoretic kernels and their existence arguments.
Composite hypotheses, unknown signal laws, estimation from data, sequential
stopping and adversarial distributional changes require their specified
models and guarantees.
\end{proofboundary}
